\documentclass[11pt]{article}
\usepackage[margin=25mm]{geometry}
\usepackage{fontspec}
\setmainfont{Latin Modern Roman}
\usepackage{amsmath,amssymb,mathtools}
\usepackage{unicode-math}
\setmathfont{Latin Modern Math}
\usepackage{xurl,hyperref}
\hypersetup{colorlinks=true,urlcolor=blue}
\setcounter{secnumdepth}{0}
\setlength{\parindent}{0pt}
\setlength{\parskip}{5pt}
\setlength{\emergencystretch}{3em}
\newcommand{\sref}[2]{\hyperlink{source:#1:#2}{[S#2]}}
\newcommand{\cataloguescope}[1]{\paragraph{Recorded target.}#1}
\begin{document}
\section{Kirby Problem 3.15}\label{statement}

\subsection{Definitions and mathematical statement}

\paragraph{Shared notation.}$\mathbb N=\{1,2,\ldots\}$, and $\mathbb Z,\mathbb Q,\mathbb R,\mathbb C$ denote the integers, rationals, reals and complex numbers. Any other conventions are specified locally.


Let $M$ be a closed oriented hyperbolic $3$-manifold, with first Betti number $b_1(M)=\dim_{\mathbb R}H^1(M;\mathbb R)>0$. Its hyperbolic metric has sectional curvature $-1$. For an oriented real rank-two bundle $E\to M$, write $e(E)\in H^2(M;\mathbb Z)$ for its Euler class and $e(E)_{\mathbb R}$ for its image in real cohomology.
A cooriented codimension-one foliation $\mathcal F$ is taut if every leaf meets a closed transverse curve; its Euler class is $e(T\mathcal F)$. A positive cooriented contact structure is a plane bundle $\xi=\ker\lambda$ with $\lambda\wedge d\lambda>0$, oriented by $d\lambda|_\xi$. It is tight if it has no overtwisted disk, and universally tight if its lift to the universal cover is tight. For an oriented nonsingular flow generated by $X$, use the Euler class of its oriented normal bundle $TM/\mathbb RX$. A pseudo-Anosov flow has transverse invariant stable and unstable two-dimensional foliations, with finitely many possible pronged singular periodic orbits, and contraction/expansion away from them. A quasigeodesic flow has constants $K\ge1,C\ge0$ such that every lifted orbit $\gamma$ in hyperbolic space, parametrized by arclength, satisfies
\[
 K^{-1}|s-t|-C\le d(\gamma(s),\gamma(t))\le K|s-t|+C.
\]
For $\rho:\pi_1(M)\to\operatorname{Homeo}^+(S^1)$, form the oriented flat circle bundle $E_\rho=(\widetilde M\times S^1)/\pi_1(M)$ using the diagonal action; its Euler class is $e(E_\rho)$.
\paragraph{Statement recorded in the supplied source.}
For each such $M$, determine the following six subsets of $H^2(M;\mathbb R)$:
\[
\begin{aligned}
\mathcal E_{\rm taut}&=\{e(T\mathcal F)_{\mathbb R}:\mathcal F\text{ taut}\},\\
\mathcal E_{\rm tight}&=\{e(\xi)_{\mathbb R}:\xi\text{ tight}\},\\
\mathcal E_{\rm universal}&=\{e(\xi)_{\mathbb R}:\xi\text{ universally tight}\},\\
\mathcal E_{\rm pA}&=\{e(TM/\mathbb RX)_{\mathbb R}:X\text{ generates a pseudo-Anosov flow}\},\\
\mathcal E_{\rm qg}&=\{e(TM/\mathbb RX)_{\mathbb R}:X\text{ generates a quasigeodesic flow}\},\\
\mathcal E_{S^1}&=\{e(E_\rho)_{\mathbb R}:\rho:\pi_1(M)\to\operatorname{Homeo}^+(S^1)\}.
\end{aligned}
\]
The question is to characterize all realizable classes, not only classes on the boundary of the dual Thurston norm ball. \sref{statement}{2}
\subsection{Short English statement}
Which real cohomology classes occur as Euler classes of taut foliations, tight or universally tight contact structures, pseudo-Anosov or quasigeodesic flows, and oriented circle actions on a closed hyperbolic three-manifold with positive first Betti number?
\subsection{Sources}
\begin{description}
\item[\hypertarget{source:statement:1}{S1}] ULAM.AI and the UnsolvedMath Contributors, \emph{UnsolvedMath: A Curated Collection of Open Mathematics Problems}, catalogue entry. CC BY 4.0. \url{https://huggingface.co/datasets/ulamai/UnsolvedMath}.
\item[\hypertarget{source:statement:2}{S2}] R. İnanç Baykur, Robion C. Kirby and Daniel Ruberman (editors), \emph{K3: A New Problem List in Low-Dimensional Topology}, Mathematical Surveys and Monographs 295 (2026), Problem 3.15, p. 141 and following remarks; original author version. \url{https://math.berkeley.edu/sites/default/files/surv-295-ruberman-watermarked-author-pdf.pdf}.
\end{description}
\label{end:statement}
\end{document}
